(Atividade \emph{Warmup C++}) Elabore um programa em C++ que receba uma sequência de valores inteiros. A sequência termina com a entrada -1. Em seguida, o programa recebe um valor $k$.

A saída do programa deve ser os $k$ menores elementos da sequência ordenados. É fortemente recomendado que se utilize os recursos da biblioteca padrão do C++ \emph{Standard Template Library} (STL) (ex. \texttt{vector}, \texttt{pair}, etc.). Observação: uma solução trivial resolve o problema em $O(nk)$ e a solução utilizando uma estrutura de dados adequada resolve em $O(klogn)$.

\inout

\begin{table}[!htb]
	\centering
	\begin{minipage}[t]{0.32\textwidth} % Primeira tabela (ajuste para 1/3 da largura)
		\centering
		\vspace{0pt}
		\begin{tabular}{cc}
			\hline
			Entrada     & Saída \\ \hline
			3 4 5 1 -1    & 1 3    \\ 
			2        &          \\
			\hline
		\end{tabular}
	\end{minipage}%
	\hfill
	\begin{minipage}[t]{0.32\textwidth} % Segunda tabela
		\centering
		\vspace{0pt}
		\begin{tabular}{cc}
			\hline
			Entrada     & Saída \\ \hline
			8 2 5 4 15 4 4 -1    & 2 4 4 4    \\ 
			4        &          \\
			\hline
		\end{tabular}
	\end{minipage}%
	\hfill
	\begin{minipage}[t]{0.32\textwidth} % Terceira tabela
		\centering
		\vspace{0pt}
		\begin{tabular}{cc}
			\hline
			Entrada    & Saída \\ \hline
			1 2 3 4 5 -1        &    1 2 3   \\
			3        &       \\ \hline
		\end{tabular}
	\end{minipage}
\end{table}